% GENERATED FILE. Edit canonical lessons, then run npm run generate:books.
% canonical-source-hash: fnv1a64:1f07b05c733a4a62
% canonical-lessons: RU-W259-with-a-clock

\chapter{Writing --- With A Clock}
\label{ch:ru-with-a-clock}
\hlchaptermodality{259}

\begin{chapteropening}
\textbf{By the end of this chapter:} \emph{I can write lines I own to a named reader inside a three-minute practice clock, and say why, under this book's 66\%-everywhere target, the weakest TORFL subtest decides.}

Complete RU-W259-with-a-clock: ten minutes on the book's clock, stop on time, then check the handwritten shapes before counting anything.
\end{chapteropening}

\section[(with a clock)]{The same writing, with a clock}

\label{lesson:RU-W259-with-a-clock}



\begin{quote}
You wrote your own line with as long as you liked.

A task you can only do untimed is not yet a task you can do.
\end{quote}

\subsection*{Writing — timed production}
\textbf{What the target publishes.} Russian's A1 exam is the state test TORFL at its Elementary Level, \textbf{\ru{ТЭУ}}. The Pushkin Institute describes it as \textbf{five separate subtests} — Lexis and Grammar, Reading, Writing, Listening and Speaking — taking about three and a half hours in all.

What this book does \textbf{not} quote is the writing paper's own minutes or length. It has not yet transcribed the current paper into \texttt{task-shapes/a1.json}, so check the provider's current materials before the exam you book.

So the clock here is this book's, not the exam's, and it is short. Set a \textbf{three-minute timer}. Write by hand to one reader you name: who you are, what you do, what you like. Lines you own, a dot at the end of each.

When the three-minute timer rings, \textbf{stop}, even mid-word.

\begin{tcolorbox}[breakable,colback=teal!4,colframe=teal!35!black,title={Before you move on}]
Spend at most one minute checking, and do not rewrite the timed lines. Hurry shows in the hand first: check \textbf{\ru{л}} against \textbf{\ru{м}}, and the bar over \textbf{\ru{т}}, before you check anything else.

Now the rule worth carrying out of this lesson, because it changes how you study rather than how you write today.

TORFL scores each subtest \textbf{on its own}, and this book's target is \textbf{66 percent in every one of the five}. Some TORFL materials describe a narrower allowance for one subtest; this book does not plan around it.

So the paper that decides your result is \textbf{the one you are worst at}, and there are five of them, not four. Strong reading cannot carry weak writing.
\end{tcolorbox}
